\documentclass[11pt]{article}
\usepackage[utf8]{inputenc}
\usepackage[T1]{fontenc}
\usepackage{lmodern}
\usepackage{amsmath,amssymb,amsthm}
\usepackage{microtype}
\usepackage[letterpaper,margin=1.1in]{geometry}
\usepackage{booktabs}
\usepackage[hidelinks]{hyperref}
\usepackage{cleveref}

\newcommand{\fhat}{\hat{f}}
\newcommand{\ghat}{\hat{g}}
\newcommand{\enc}{\operatorname{enc}}
\newcommand{\dec}{\operatorname{dec}}
\newcommand{\B}{\{0,1\}}
\newcommand{\TV}{\operatorname{TV}}
\newcommand{\im}{\operatorname{im}}
\newcommand{\orbitF}{\operatorname{orbit}_F}

\theoremstyle{definition}
\newtheorem{definition}{Definition}
\theoremstyle{plain}
\newtheorem{proposition}{Proposition}
\theoremstyle{remark}
\newtheorem{remark}{Remark}

\crefname{definition}{Definition}{Definitions}

\title{A Quantitative Security Model for Trapdoor Computing}
\author{Alexander Towell\\\texttt{lex@metafunctor.com}}
\date{2026}

\begin{document}
\maketitle

\begin{abstract}
Trapdoor computing lets an untrusted machine evaluate opaque bit strings
through an opaque lookup table while only a trusted machine encodes and
decodes. Its central claim is that confidentiality is \emph{measured}, an
explicit leakage quantity a designer can trade, rather than \emph{negligible}
in a security parameter. That claim has so far rested on an informal model: the
constructions come with leakage \emph{measures}, but not with a security
\emph{definition} tying an adversary's power to what is provably guaranteed.
This paper supplies the missing model. We formalize the trusted/untrusted
adversary and its view, and define the confidentiality of a cipher map in the
language of quantitative information flow, as a tuple of bounded, composable
leakage quantities rather than a negligible advantage. We show that the
framework's three leakage measures, the marginal entropy ratio, the
compositional joint-recovery rate, and the active orbit-closure bound, are
three aspects of one adversary with three kinds of access, and that the
resulting security tuple composes under cipher-map composition. We close by
relating the definition to simulation-based and game-based cryptographic
security, arguing that quantitative information flow, not indistinguishability,
is the right model for the measurable-leakage regime trapdoor computing
occupies.
\end{abstract}

\section{Introduction}
\label{sec:intro}

A security definition does two things: it says what an adversary may do, and it
says what is guaranteed against an adversary who does it. Trapdoor computing has
a clear adversary, the untrusted machine, and a clear collection of guarantees,
the leakage bounds its constructions satisfy, but it has not had a definition
that joins the two. The cipher-maps abstraction states informally that the
untrusted machine ``cannot decode, cannot distinguish a real query from filler,
cannot determine the domain, cannot be sure a result is correct''~%
\cite{towell2026cipher}; the confidentiality theory gives an entropy-ratio
bound on a single token and a recovery-rate bound on a chain~%
\cite{towell2026maxconf}; the type theory gives an orbit-closure bound on an
active adversary~\cite{towell2026algebraic}. These are measures in search of a
model.

The gap is not cosmetic. A reviewer trained on simulation-based security asks,
reasonably, for a definition: an adversary, a security experiment, and a
statement of what the experiment guarantees. Absent one, ``measured, not
negligible'' reads as a slogan rather than a claim. This paper makes it a claim.
We give the formal trusted/untrusted model and the adversary's view
(\Cref{sec:model}); define a cipher map's confidentiality as a tuple of bounded
leakage quantities in the sense of quantitative information flow~%
\cite{smith2009foundations,alvim2020science} (\Cref{sec:definition}); show the
framework's three leakage measures are one adversary at three access levels
(\Cref{sec:scales}); prove the security tuple composes (\Cref{sec:composition});
and defend the choice of a quantitative rather than a simulation-based or
game-based definition (\Cref{sec:crypto}).

The contribution is the definition itself, and the demonstration that the
framework's existing, separately-proven bounds are instances of it. We prove no
new leakage bound here; we give the model the bounds were always bounds in.

\section{The Trusted/Untrusted Model}
\label{sec:model}

Fix a latent function $f : X \to Y$ known only to a trusted party.

\begin{definition}[Cipher map; from~\cite{towell2026cipher}]
\label{def:cipher-map}
A \emph{cipher map} for $f$ is a tuple $(\fhat, \enc, \dec, s)$ with
$\fhat : \B^n \to \B^n$ a total function, $\enc : X \times \{0, \ldots,
K(x)-1\} \to \B^n$, $\dec : \B^n \to Y \cup \{\bot\}$, and $s$ a secret from
which all three derive. The \emph{trusted machine} $T$ holds $(\enc, \dec, s)$;
the \emph{untrusted machine} $U$ holds $\fhat$.
\end{definition}

The adversary is $U$. We make its power and its view explicit.

\begin{definition}[Adversary view]
\label{def:view}
In a run with in-domain inputs $x_1, \ldots, x_N$ drawn i.i.d.\ from a query
distribution $D$ on $X$, the untrusted machine's view is
\[
  \mathsf{view}_U \;=\; \bigl(\, \fhat,\; (c_i)_{i=1}^N,\; (\fhat(c_i))_{i=1}^N,\;
  F \,\bigr),
\]
where $c_i = \enc(x_i, k_i)$ is the $i$-th query token, and $F$ is the set of
cipher maps $U$ holds and may apply to any token to produce and observe further
tokens. $U$ is unbounded for the purpose of the leakage analysis: we measure how
much $\mathsf{view}_U$ reveals, not whether a bounded $U$ can extract it.
\end{definition}

Two modeling choices deserve comment. First, $\fhat$ is \emph{public to $U$}: the
untrusted machine holds the whole table, so the model grants the adversary the
construction itself and asks what the construction still hides. Second, a
\emph{malicious} $U$ that deviates from the protocol gains nothing the model
does not already grant, because totality makes every token, honest or not,
produce an answer; there is no error channel to probe. The honest-but-curious
$U$ of \Cref{def:view} is therefore the worst case for leakage.

The latent quantities $U$ would learn, and that the model must bound, are the
inputs $X = (x_i)$, the function $f$, the distribution $D$, and quantities
derived from them: the domain, the equality structure of the queries, and the
joint distribution of correlated values across several cipher maps.

\section{Confidentiality as Bounded Leakage}
\label{sec:definition}

Quantitative information flow measures the leakage of a secret $S$ to a view
$V$ by how much $V$ reduces uncertainty about $S$, for instance the Shannon
leakage $I(S; V)$ or a min-entropy (Bayes-vulnerability) leakage. Security in
this language is not the absence of leakage but a \emph{bound} on it. We define
a cipher map's confidentiality as a tuple of such bounds, one per kind of
access in \Cref{def:view}.

\begin{definition}[Confidentiality of a cipher map]
\label{def:confidentiality}
Let $(\fhat, \enc, \dec, s)$ be a cipher map with representation uniformity
$\delta = \TV(Q, U_{\im})$ on its populated support, where $Q$ is the induced
token distribution and $U_{\im}$ is uniform on $\im(\enc)$, and let $H^{*} =
\log_2|\im(\enc)|$. The cipher map is \emph{$(\delta, R, F)$-confidential}
against the adversary of \Cref{def:view} if all three of the following hold.
\begin{description}
\item[\textmd{(M) Marginal.}] For a single token $c = \enc(x, k)$ with
  $x \sim D$, the leakage of $x$ is bounded by the entropy ratio:
  \[
    \frac{I(x; c)}{H^{*}} \;\le\; 1 - e, \qquad
    e = \frac{H(Q)}{H^{*}} \;\ge\; 1 - \delta - \frac{h_2(\delta)}{H^{*}},
  \]
  the lower bound on $e$ being the Fannes--Audenaert inequality, linear in
  $\delta$. Small $\delta$ bounds the marginal leakage.
\item[\textmd{(C) Compositional.}] For two cipher maps $\fhat_1, \fhat_2$ for
  $f_1 : X \to Y_1$, $f_2 : X \to Y_2$ observed on shared tokens, mutual
  information is preserved, $I(\fhat_1(C); \fhat_2(C)) = I(f_1(X); f_2(X))$, and
  the latent joint distribution on $Y_1 \times Y_2$ is recoverable to
  total-variation accuracy $\xi$ from $N = R(\xi)$ observations with
  $R(\xi) = \Theta(|Y_1|\,|Y_2| / \xi^2)$.
\item[\textmd{(A) Active.}] For an adversary applying the operations $F$ to a
  token $c$, the residual entropy of the latent value $X$ behind $c$ is
  \[
    H(X \mid \mathsf{view}_U) \;\ge\; H(X) - \log_2|\orbitF(c)|,
  \]
  where $\orbitF(c)$ is the set of tokens reachable from $c$ under $F$.
\end{description}
\end{definition}

The definition is parameterized two ways: by the \emph{construction} (through
$\delta$ and the acceptance partition that shapes the orbit and the recovery
constants) and by the \emph{adversary's access} (through the number of shared
observations $N$ and the operation set $F$). A cipher map is not ``secure'' or
``insecure'' absolutely; it is secure to these explicit, computable degrees, and
the degrees are exactly the quantities a designer trades against space,
bandwidth, and accuracy. This is the formal content of ``measured, not
negligible.''

\begin{remark}[Why a tuple, not a single number]
\label{rem:tuple}
The three clauses cannot be collapsed into one bound, because they are not
comparable quantities. (M) is a normalized leakage of a single value, governed
by a construction parameter. (C) is a \emph{rate} at which a joint distribution
becomes estimable, governed by the latent cardinalities and unaffected by any
construction parameter. (A) is an unnormalized residual entropy, governed by the
exposed operations. A security model for trapdoor computing must carry all three
because a cipher map can be strong on one and weak on another: near-perfect
marginal confidentiality ($\delta \approx 0$) coexists with full compositional
recovery given enough shared observations. The irreducibility of the scales is
a feature of the model, not a defect of a construction.
\end{remark}

\section{Three Scales, One Adversary}
\label{sec:scales}

The three clauses of \Cref{def:confidentiality} are often read as three
adversaries; they are one adversary at three levels of access (\Cref{tab:scales}).
Clause (M) is the adversary of \Cref{def:view} restricted to a single token and
passive. Clause (C) is the same adversary, still passive, but observing several
cipher maps on shared tokens, the access that exposes joint structure. Clause
(A) is the adversary using the last component of its view, the operations $F$,
to probe actively. Nothing in the three requires different powers; they
partition the single view of \Cref{def:view} by which part of it the adversary
exploits.

\begin{table}[t]
\centering
\caption{The three clauses are one adversary (\Cref{def:view}) at three access
levels.}
\label{tab:scales}
\begin{tabular}{@{}llll@{}}
\toprule
Clause & Access used & Governed by & Bounded by $\delta$? \\
\midrule
(M) marginal & one token, passive & $\delta$ (entropy ratio) & yes \\
(C) compositional & shared tokens, passive & $|Y_1||Y_2|$, $N$ & no \\
(A) active & operations $F$ & $|\orbitF(c)|$ & no \\
\bottomrule
\end{tabular}
\end{table}

The table makes the model's sharpest statement visible. Only (M) is bought by a
construction parameter. The compositional and active leakages are governed by
quantities $\delta$ does not touch, the latent cardinalities and the number of
observations for (C), the reachable orbit for (A). A designer who lowers
$\delta$ to its floor has done everything the marginal scale allows and nothing
for the other two; defending those is a matter of system design, limiting shared
observations and exposed operations, not of tuning a parameter. The model thus
predicts where engineering effort can and cannot help.

\section{Composition}
\label{sec:composition}

A security definition is useful for building systems only if it composes: the
security of a composite must be computable from the security of its parts. The
tuple of \Cref{def:confidentiality} does.

\begin{proposition}[Composition of the security tuple]
\label{prop:composition}
Let $\fhat_1, \ldots, \fhat_m$ be cipher maps chained so each consumes the
previous output, with correctness parameters $\eta_i$ and marginal bounds from
$\delta_i$. Then the chain is a cipher map whose security tuple is determined by
the parts:
\begin{enumerate}
\item correctness composes, $\eta_{\mathrm{total}} \le 1 - \prod_i (1 - \eta_i)$;
\item the marginal leakage of the chain's observable token stream is bounded by
  the entropy ratio of that stream, which the per-stage $\delta_i$ bound;
\item the compositional rate and the active orbit of the chain grow predictably,
  the orbit by at most the product of the per-step frontiers and the recovery
  rate by the latent cardinalities the chain exposes.
\end{enumerate}
\end{proposition}

The correctness bound is the elementary observation that a chain is correct when
every stage is, an inequality rather than an equality because a stage error can
be masked downstream~\cite{towell2026cipher}. The marginal and active parts
inherit their composition from the entropy-ratio and orbit-closure machinery the
framework already proves. The point for the model is structural: because each
clause composes, a system assembled from cipher maps has a security tuple a
designer can compute from a parts list, which is the property that lets the
definition support construction and not only analysis.

\section{Relation to Cryptographic Security}
\label{sec:crypto}

The natural objection to \Cref{def:confidentiality} is that it is not a
cryptographic security definition. It is not, and the difference is the point.

\paragraph{Against simulation-based security.}
Searchable encryption and oblivious RAM are proved secure by exhibiting a
simulator that reproduces the adversary's view from a \emph{leakage function}~%
\cite{curtmola2006searchable,goldreich1996software}: security means the real
view is indistinguishable from a simulation that depends only on what the
leakage function reveals. Trapdoor computing's leakage is not a function to be
fed to a simulator; it is a family of \emph{bounds}, and the constructions do
not claim that the residual is simulatable, only that it is small in the
measured sense. One could define a leakage function and seek a simulator, but
the resulting statement would say less than \Cref{def:confidentiality}: a
simulator with a large leakage function certifies indistinguishability while
hiding how much leaks, whereas the quantitative bound says exactly how much.

\paragraph{Against game-based security.}
Indistinguishability definitions promise a \emph{negligible} advantage in a
security parameter. Trapdoor computing makes no such promise, and a construction
forced to claim one would either fail, because its compositional leakage is not
negligible, or mislead. The honest statement is the quantitative one: the
advantage is a measured quantity, often not negligible, and useful precisely
because it is known.

\paragraph{Why quantitative information flow is the right model.}
The leakage trapdoor computing incurs is operationally meaningful (its bounds
are the quantities a designer trades), composable (\Cref{prop:composition}), and
honest (it does not assert a negligibility it cannot support). These are exactly
the virtues quantitative information flow was built to express. The cost is real:
the guarantee is weaker than indistinguishability, and for an application that
can afford the strong tools, the strong tools are correct. For the regime
between plaintext and oblivious RAM, where leakage is unavoidable but bounded
and tradeable, the quantitative definition is not a weaker substitute for the
cryptographic one; it is the right definition.

\section{Conclusion}

We have given trapdoor computing the security model its measurable-leakage claim
required: a formal adversary and view, a definition of confidentiality as a
tuple of bounded leakage quantities in the sense of quantitative information
flow, a demonstration that the framework's three leakage measures are one
adversary at three access levels, and a proof that the tuple composes. The model
proves no new bound; it gives the bounds the framework already had a definition
to be bounds in, and in doing so answers the question a foundations reviewer asks
first: what does it mean, here, to be secure? It means leaking a measured,
composable, and computable amount, and no more.

\paragraph{Provenance.} The trapdoor framework was begun in 2017 and developed
through 2020, then recovered and extended in 2024; this model formalizes the
security stance implicit in it throughout.

\bibliographystyle{plain}
\bibliography{references}

\end{document}
